\section{Liquid Templates}
\label{app:liquid-templates}

\href{https://shopify.github.io/liquid}{\sffamily Liquid} is a lightweight,
text-oriented templating language designed to separate \emph{content} from
\emph{presentation}. %
Rather than embedding data or logic directly in a document, Liquid allows
the writer to design a template that describes the structure of the final
output, while delegating the actual content to external data sources such
as \seesection{app:yaml}{\YAML} or \JSON{} files. %
During rendering, a Liquid engine combines the template and the data to
produce a complete, fully expanded document. %

\medskip

A Liquid \emph{template} consists of ordinary text enriched with a set of
tag types. %
These tags make it possible to insert values, perform conditional
operations, iterate over collections, and apply simple transformations. %

\begin{itemize}
  \item \textbf{Output tags} (e.g.\ \liq{{{ title }}}) insert
        the value of a variable directly into the document.
  \item \textbf{Logic tags} enable conditional inclusion of
    content, e.g.:
\begin{lqdblock}
{% if x > 0 %}
 ...
{% endif %}
\end{lqdblock}
  \item \textbf{Iteration tags} repeat a block of text for
    each element of a sequence, e.g.:
\begin{lqdblock}
{% for book in books %}
  {{ book.title }} ({{ book.author }})
{% endfor %}
\end{lqdblock}
  \item \textbf{Filters} (e.g.\ \liq{{{ name | upcase }}}) apply
        simple transformations to values before insertion.
\end{itemize}

\medskip

Liquid is of course a lot more powerful than these examples, and the
interested reader will find a full reference of the templating system here:
\begin{center}
  \url{https://shopify.github.io/liquid/}
\end{center}

%%% Local Variables:
%%% mode: LaTeX
%%% TeX-master: "../liquidtex"
%%% End:
